% trienerments.tex — Formal writeup on t(r)ienerments
% oracle-2026-05-02-S1
\documentclass[12pt,reqno]{amsart}
\usepackage[margin=1.2in,top=1.1in,bottom=1.2in]{geometry}
\usepackage[T1]{fontenc}
\usepackage[utf8]{inputenc}
\usepackage[protrusion=true,expansion=false]{microtype}
\usepackage{amsmath,amssymb,amsthm,mathtools}
\usepackage{enumitem}
\usepackage{booktabs,array}
\usepackage{xcolor}
\usepackage{hyperref}
\hypersetup{colorlinks=true,linkcolor=blue,citecolor=blue,urlcolor=blue}
\usepackage{tikz}
\usetikzlibrary{arrows.meta,positioning,calc}

% ── Theorem environments ────────────────────────────────────────────────────
\newtheorem{theorem}{Theorem}[section]
\newtheorem{lemma}[theorem]{Lemma}
\newtheorem{proposition}[theorem]{Proposition}
\newtheorem{corollary}[theorem]{Corollary}
\theoremstyle{definition}
\newtheorem{definition}[theorem]{Definition}
\newtheorem{example}[theorem]{Example}
\newtheorem{remark}[theorem]{Remark}
\theoremstyle{plain}
\newtheorem{conjecture}[theorem]{Conjecture}

% ── Macros ───────────────────────────────────────────────────────────────────
\newcommand{\Trie}{\mathcal{T}}          % a t(r)ienerment
\newcommand{\Trien}[1]{\mathcal{T}_{#1}} % specific t(r)ienerment
\newcommand{\Sn}[1]{S_{#1}}             % symmetric group
\newcommand{\Knv}[1]{K_{#1}}            % complete graph
\newcommand{\HP}[1]{H(#1)}              % Hamiltonian path count
\newcommand{\Fix}{\mathrm{Fix}}
\newcommand{\lcm}{\mathrm{lcm}}
\newcommand{\floor}[1]{\lfloor #1 \rfloor}
\newcommand{\Omega}{\Omega}
\newcommand{\IsoClass}[1]{[#1]}
\newcommand{\OEIS}[1]{\href{https://oeis.org/#1}{\texttt{#1}}}

\allowdisplaybreaks

% ── Title ────────────────────────────────────────────────────────────────────
\title{T(r)ienerments: Ternary Generalizations of Tournaments\\
and Their Isomorphism Theory}

\author{Eliott Cassidy}
\address{Multi-Agent Mathematics Research Project}
\email{guideposttutors@gmail.com}
\date{May 2, 2026}

\begin{document}
\maketitle

\begin{abstract}
A \emph{t(r)ienerment} is a complete graph on $n$ vertices in which each edge is
assigned one of three states: a directed arc $u \to v$, a directed arc $v \to u$,
or a bidirectional \emph{tie} $u \leftrightarrow v$.  Ties allow Hamiltonian paths
to be formed through them in either direction, making t(r)ienerments a natural
ternary generalization of tournaments.

We derive a closed Burnside formula for the number of unlabeled isomorphism classes,
which equals the sequence \OEIS{A007025} of unlabeled oriented graphs.  We establish
a canonical isomorphism between ``positive'' t(r)ienerments (ties as $\leftrightarrow$)
and ``negative'' t(r)ienerments (ties as absent edges), explaining why both spaces
are counted by the same sequence.

We introduce and compute the refined distribution $f(n,k)$, the number of
isomorphism classes on $n$ vertices with exactly $k$ ties.  The $k=0$ slice
recovers the tournament count \OEIS{A000568}, while $f(n,\binom{n}{2})=
f\!\left(n,\binom{n}{2}-1\right)=1$ for all $n$.  The tail values stabilize:
$f\!\left(n,\binom{n}{2}-2\right)=4$ for $n\ge 4$, and
$f\!\left(n,\binom{n}{2}-3\right)=14$ for $n\ge 6$.

We prove that every t(r)ienerment has at least one directed Hamiltonian path
(extending R\'edei's theorem), and show that, unlike tournaments, the path
count $H(\Trie)$ need not be odd.  The Grinberg--Stanley Odd-Cycle
Collection Formula $H(\Trie)=I(\Omega(D_\Trie),2)$ extends to all
t(r)ienerments via the associated digraph $D_\Trie$.
\end{abstract}

\tableofcontents

% ════════════════════════════════════════════════════════════════════════════
\section{Introduction}
% ════════════════════════════════════════════════════════════════════════════

A \emph{tournament} on $n$ vertices is a complete directed graph: for every pair
$\{u,v\}$, exactly one of $u\to v$ or $v\to u$ holds.  Tournaments model
ranking situations (round-robin competitions, paired comparisons, elections) in
which every match produces a strict winner.

Real-world ranking data, however, frequently involves \emph{ties}: two
alternatives that cannot be separated by the comparison mechanism.  A referee
may declare a draw; two items may score identically on a metric; rounding in
a sports table may leave two teams indistinguishable.  The natural question
is how to extend tournament theory to accommodate a third outcome.

\subsection{Two representations of ties}

There are two canonical ways to incorporate ties into a complete graph on $n$
vertices.

\begin{itemize}
\item \textbf{Positive representation (t(r)ienerments).}
  Each pair $\{u,v\}$ receives one of three labels: $u\to v$, $v\to u$, or
  $u\leftrightarrow v$.  Bidirectional edges are the ``positive'' encoding of a
  tie: they exist in both directions and may be traversed freely by a
  Hamiltonian path.

\item \textbf{Negative representation (oriented graphs).}
  Each pair $\{u,v\}$ receives one of three labels: $u\to v$, $v\to u$, or
  $\emptyset$ (no arc).  Absent edges encode ties: they represent outcomes that
  were never decided.  These objects are precisely the \emph{oriented graphs}
  (simple digraphs with no symmetric pair of arcs) studied in graph theory.
\end{itemize}

A central observation of this paper (Theorem~\ref{thm:positive-negative-iso})
is that these two representations give \emph{isomorphic} combinatorial spaces:
the $S_n$-equivariant bijection that swaps ties and absent edges preserves all
structure, so the two spaces have identical isomorphism-class counts.

\subsection{The name}

We call an object in the positive representation a \emph{t(r)ienerment},
blending ``tournament'' with the prefix ``tri'' (three outcomes) and a nod to
the Wiener score used in some tournament-ranking contexts.  The parenthesised
``r'' signals both ``ternary'' and the alternative reading ``trienerment''
(three-outcome tournament).

\subsection{Summary of results}

\begin{enumerate}
\item The number of unlabeled isomorphism classes of $n$-vertex t(r)ienerments
  is
  \[
    T(n)
    = \frac{1}{n!}\sum_{\sigma\in S_n} 3^{B(\sigma)},
    \quad
    B(\sigma)=\sum_a\floor{\tfrac{l_a-1}{2}}+\sum_{a<b}\gcd(l_a,l_b),
  \]
  where $(l_1,l_2,\ldots)$ is the cycle type of $\sigma$.
  The sequence $T(n)$ is \OEIS{A007025}.

\item The refined distribution $f(n,k)$ (iso classes with $k$ ties)
  is computed from a generating polynomial per cycle type,
  with $f(n,0)=\text{A000568}(n)$ and
  $f(n,\binom{n}{2})=f(n,\binom{n}{2}-1)=1$.

\item Every t(r)ienerment has at least one directed Hamiltonian path
  (R\'edei's theorem extends), but $H(\Trie)$ need not be odd.

\item The Odd-Cycle Collection Formula extends:
  $H(\Trie)=I(\Omega(D_\Trie),2)$ for the associated digraph $D_\Trie$.
\end{enumerate}

\subsection{Relation to existing work}

Tournament isomorphism classes are counted by the Burnside formula
\[
  \text{A000568}(n)=\frac{1}{n!}\sum_{\substack{\sigma\in S_n\\
    \text{all cycles odd}}} 2^{B(\sigma)},
\]
where the restriction to odd-cycle permutations arises because even cycles
force a symmetric pair of arcs, which is forbidden in tournaments.
In t(r)ienerments, bidirectional ties resolve this obstruction, so all of
$S_n$ contributes with weight $3^{B(\sigma)}$.

The sequence $T(n)=\text{A007025}(n)$ counts unlabeled oriented graphs, which
appear in the context of ``semi-complete digraphs'' and ``tournaments with
draws'' in social choice theory.  The connection to the triangular distribution
$f(n,k)$ appears to be new.

\subsection{Organization}

Section~\ref{sec:defs} sets up formal definitions.  Section~\ref{sec:burnside}
derives the Burnside formula.  Section~\ref{sec:distribution} develops the
refined distribution $f(n,k)$.  Section~\ref{sec:positive-negative} proves the
positive-negative isomorphism.  Section~\ref{sec:hamiltonian} discusses
Hamiltonian paths and the OCF extension.  Section~\ref{sec:tables} collects
the computed tables.  Section~\ref{sec:open} lists open questions.

% ════════════════════════════════════════════════════════════════════════════
\section{Definitions}
\label{sec:defs}
% ════════════════════════════════════════════════════════════════════════════

\begin{definition}[T(r)ienerment]\label{def:trienerment}
Let $V=\{1,\ldots,n\}$.  A \emph{t(r)ienerment} on $V$ is a function
\[
  \tau: \binom{V}{2}\to \bigl\{\text{``}u\to v\text{''}, \text{``}v\to u\text{''}, \text{``}u\leftrightarrow v\text{''}\bigr\}
\]
that assigns to each unordered pair $\{u,v\}$ one of three states.  We write:
\begin{itemize}
\item $u\xrightarrow{\tau}v$ if $\tau(\{u,v\})=\text{``}u\to v\text{''}$,
\item $v\xrightarrow{\tau}u$ if $\tau(\{u,v\})=\text{``}v\to u\text{''}$,
\item $u\xleftrightarrow{\tau}v$ if $\tau(\{u,v\})=\text{``}u\leftrightarrow v\text{''}$.
\end{itemize}
The number of \emph{ties} in $\tau$ is
$k(\tau)=\bigl|\bigl\{\{u,v\}:\tau(\{u,v\})=\text{``}u\leftrightarrow v\text{''}\bigr\}\bigr|$.
\end{definition}

\begin{remark}
When $k(\tau)=0$, a t(r)ienerment is exactly a tournament.
\end{remark}

\begin{definition}[Isomorphism]\label{def:iso}
Two t(r)ienerments $\tau,\tau'$ on $V$ are \emph{isomorphic}, written $\tau\cong\tau'$,
if there exists a permutation $\sigma\in S_n$ such that for all $\{u,v\}\in\binom{V}{2}$,
$\tau'(\{\sigma(u),\sigma(v)\})$ is the same state as $\tau(\{u,v\})$ with $u$ and $v$
replaced by $\sigma(u)$ and $\sigma(v)$.  (In particular, the state ``$u\to v$'' maps to
``$\sigma(u)\to\sigma(v)$'', and the state ``$u\leftrightarrow v$'' maps to
``$\sigma(u)\leftrightarrow\sigma(v)$''.)
\end{definition}

\begin{definition}[Associated digraph]\label{def:digraph}
The \emph{associated digraph} $D_\tau$ of a t(r)ienerment $\tau$ is the digraph
on $V$ with arc set
\[
  A(D_\tau)=\bigl\{(u,v):u\xrightarrow{\tau}v\text{ or }u\xleftrightarrow{\tau}v\bigr\}.
\]
Each directed edge $u\to v$ in $\tau$ contributes the single arc $(u,v)$;
each tie $u\leftrightarrow v$ contributes both arcs $(u,v)$ and $(v,u)$.
\end{definition}

\begin{definition}[Directed Hamiltonian path]
A \emph{directed Hamiltonian path} of a t(r)ienerment $\tau$ is a sequence
$v_1,v_2,\ldots,v_n$ of all $n$ vertices such that $(v_i,v_{i+1})\in A(D_\tau)$
for $i=1,\ldots,n-1$.  Equivalently, for each consecutive pair $(v_i,v_{i+1})$,
either $v_i\xrightarrow{\tau}v_{i+1}$ or $v_i\xleftrightarrow{\tau}v_{i+1}$.
We denote by $H(\tau)$ the number of directed Hamiltonian paths in $\tau$.
\end{definition}

\begin{definition}[Negative t(r)ienerment / oriented graph]
A \emph{negative t(r)ienerment} (or \emph{oriented graph}) on $V$ is a function
\[
  \nu:\binom{V}{2}\to\bigl\{\text{``}u\to v\text{''}, \text{``}v\to u\text{''}, \emptyset\bigr\}
\]
where $\emptyset$ represents a \emph{tie by absence}: no arc between $u$ and $v$.
\end{definition}

% ════════════════════════════════════════════════════════════════════════════
\section{Counting Isomorphism Classes via Burnside}
\label{sec:burnside}
% ════════════════════════════════════════════════════════════════════════════

Let $\mathfrak{T}_n$ denote the set of all labeled t(r)ienerments on $V=\{1,\ldots,n\}$.
We have $|\mathfrak{T}_n|=3^{\binom{n}{2}}$.  The group $S_n$ acts on $\mathfrak{T}_n$
by relabeling vertices: $(\sigma\cdot\tau)(\{u,v\})=\tau(\{\sigma^{-1}(u),\sigma^{-1}(v)\})$,
with the state's orientation updated accordingly.  The number of isomorphism classes
is given by Burnside's lemma:
\[
  T(n)=\frac{1}{n!}\sum_{\sigma\in S_n}|\Fix(\sigma)|,
\]
where $\Fix(\sigma)=\{\tau\in\mathfrak{T}_n:\sigma\cdot\tau=\tau\}$.

\subsection{Orbit analysis for a fixed permutation}

Fix $\sigma\in S_n$ with cycle type $(l_1,l_2,\ldots,l_m)$ (so $\sum l_a=n$).
We analyze the orbits of $\sigma$ acting on ordered pairs $(u,v)$ with $u\ne v$.

\begin{definition}
For an orbit $O$ of $\sigma$ on ordered pairs, let $O^{\mathrm{rev}}=\{(v,u):(u,v)\in O\}$.
\begin{itemize}
\item $O$ is a \emph{Case~A orbit} if $O=O^{\mathrm{rev}}$.
\item If $O\ne O^{\mathrm{rev}}$, then $(O,O^{\mathrm{rev}})$ is a \emph{Case~B orbit pair}.
\end{itemize}
\end{definition}

The orbit analysis splits into within-cycle and cross-cycle orbits.

\subsubsection{Within-cycle orbits}

For a cycle $C=(c_0,c_1,\ldots,c_{l-1})$ of length $l$, the orbit of the pair
$(c_0,c_k)$ under $\sigma$ has size $l$ and is
$O_k=\{(c_i,c_{(i+k)\bmod l}):i=0,\ldots,l-1\}$.
Its reverse is $O_{l-k}$.

\begin{lemma}\label{lem:within-cycle}
In a cycle of length $l$:
\begin{enumerate}
\item If $l$ is even, then $O_{l/2}$ is the unique Case~A orbit.  It covers
  $l/2$ unordered pairs.  There are $\floor{(l-1)/2}$ Case~B orbit pairs,
  each covering $l$ unordered pairs.
\item If $l$ is odd, there are no Case~A orbits.  There are $(l-1)/2$ Case~B
  orbit pairs, each covering $l$ unordered pairs.
\end{enumerate}
In both cases, the number of Case~B orbit pairs within a cycle of length $l$ is
$\floor{(l-1)/2}$.
\end{lemma}

\subsubsection{Cross-cycle orbits}

For two cycles $C_a$ of length $l_a$ and $C_b$ of length $l_b$ ($a\ne b$),
the orbit of a cross-pair $(c_a^i,c_b^j)$ has size $\lcm(l_a,l_b)$.

\begin{lemma}\label{lem:cross-cycle}
Between cycles $C_a$ and $C_b$, there are exactly $\gcd(l_a,l_b)$ Case~B orbit
pairs, each covering $\lcm(l_a,l_b)$ unordered pairs.  There are no Case~A
cross-cycle orbits.
\end{lemma}

\begin{proof}
The number of forward orbits from $C_a$ to $C_b$ is $l_a l_b/\lcm(l_a,l_b)=\gcd(l_a,l_b)$.
Since $C_a\ne C_b$, no cross-orbit can equal its own reverse: $(u,v)$ with $u\in C_a$,
$v\in C_b$ has $(v,u)$ in the reverse orbit which lies from $C_b$ to $C_a$, a
distinct family.
\end{proof}

\subsubsection{Consequences for fixed t(r)ienerments}

A labeled t(r)ienerment $\tau$ is fixed by $\sigma$ if and only if $\tau$ is
constant on each orbit of $\sigma$ on ordered pairs (respecting the directional
encoding: the state ``$u\to v$'' is sent to ``$\sigma(u)\to\sigma(v)$'').

\begin{lemma}\label{lem:case-a-forced}
For a Case~A orbit $O=O^{\mathrm{rev}}$ covering a set $P$ of unordered pairs:
$\tau$ is fixed by $\sigma$ on $P$ if and only if all pairs in $P$ are ties.
There is exactly one such choice.
\end{lemma}

\begin{proof}
If $(u,v)\in O$ and $(v,u)\in O$, then fixing $\tau$ requires $\tau(\{u,v\})$
to be consistent with both orderings.  The states ``$u\to v$'' and ``$v\to u$''
would need to simultaneously hold, which is impossible.  Only the symmetric
state ``$u\leftrightarrow v$'' is consistent.  Since $\tau$ is constant on $O$,
all pairs in $P$ are forced to be ties.
\end{proof}

\begin{lemma}\label{lem:case-b-free}
For a Case~B orbit pair $(O,O^{\mathrm{rev}})$ covering a set $P$ of unordered
pairs: $\tau$ is fixed by $\sigma$ on $P$ if and only if $\tau$ assigns the
same state to all pairs in $P$, and that state may be any of the three options.
There are exactly three choices.

In terms of ties: either all pairs in $P$ are directed (two choices: constant
``$u\to v$'' direction or constant ``$v\to u$'' direction, each self-consistent),
or all pairs in $P$ are ties.
\end{lemma}

\begin{proof}
Since $O$ and $O^{\mathrm{rev}}$ are distinct orbits, the values of $\tau$ on
$(u,v)$ (via $O$) and $(v,u)$ (via $O^{\mathrm{rev}}$) can be chosen independently.
Setting $f(u,v)=a\in\{0,1\}$ (where 1 means arc exists) and $f(v,u)=b\in\{0,1\}$
with $(a,b)\ne(0,0)$ (required by the t(r)ienerment constraint) gives three choices:
$(1,0)$ (directed one way), $(0,1)$ (directed the other), $(1,1)$ (tie).
\end{proof}

\begin{theorem}[Fixed-point count]\label{thm:fix}
For $\sigma\in S_n$ with cycle type $(l_1,\ldots,l_m)$,
\[
  |\Fix(\sigma)| = 3^{B(\sigma)},
  \qquad
  B(\sigma)=\sum_{a=1}^m\floor{\frac{l_a-1}{2}}+\sum_{1\le a<b\le m}\gcd(l_a,l_b).
\]
\end{theorem}

\begin{proof}
Each Case~A orbit contributes exactly 1 choice (Lemma~\ref{lem:case-a-forced}).
Each Case~B orbit pair contributes exactly 3 choices (Lemma~\ref{lem:case-b-free}).
The total number of Case~B orbit pairs is, by Lemmas~\ref{lem:within-cycle}
and~\ref{lem:cross-cycle}:
\[
  B(\sigma)
  =\underbrace{\sum_a\floor{\frac{l_a-1}{2}}}_{\text{within-cycle}}
  +\underbrace{\sum_{a<b}\gcd(l_a,l_b)}_{\text{cross-cycle}}.
  \qedhere
\]
\end{proof}

\begin{theorem}[Total isomorphism class count]\label{thm:total-count}
The number of unlabeled isomorphism classes of t(r)ienerments on $n$ vertices is
\[
  T(n)=\frac{1}{n!}\sum_{\lambda\vdash n}
  \frac{n!}{\prod_j l_j^{a_j}\cdot a_j!}\cdot 3^{B(\lambda)},
\]
where the sum is over all cycle types $\lambda=(l_1^{a_1}l_2^{a_2}\cdots)$ of $S_n$,
and $B(\lambda)$ is computed from the cycle lengths.  The values are listed in
Table~\ref{tab:T-values} and equal \OEIS{A007025}.
\end{theorem}

\begin{remark}[Comparison with tournament Burnside formula]
The tournament count \OEIS{A000568} is:
\[
  \text{A000568}(n)=\frac{1}{n!}\sum_{\substack{\sigma\in S_n\\\text{all cycles odd}}}2^{B(\sigma)}.
\]
Permutations with even cycles contribute $|\Fix(\sigma)|=0$ to tournaments (since Case~A
orbits within even cycles force symmetric arc pairs, which tournaments forbid).
T(r)ienerments allow ties, so all of $S_n$ contributes, and the base changes from 2 to 3.
Formally:
\[
  T(n)=\frac{1}{n!}\sum_{\sigma\in S_n}3^{B(\sigma)},
  \qquad
  \text{A000568}(n)=\frac{1}{n!}\sum_{\substack{\sigma\in S_n\\\text{all cycles odd}}}2^{B(\sigma)}.
\]
The t(r)ienerment formula is strictly larger and involves all permutations.
\end{remark}

% ════════════════════════════════════════════════════════════════════════════
\section{The Refined Distribution \texorpdfstring{$f(n,k)$}{f(n,k)}}
\label{sec:distribution}
% ════════════════════════════════════════════════════════════════════════════

Let $f(n,k)$ denote the number of isomorphism classes of t(r)ienerments on $n$
vertices with exactly $k$ ties.  Clearly $\sum_{k=0}^{\binom{n}{2}}f(n,k)=T(n)$.

\subsection{Generating polynomial per permutation}

To extract $f(n,k)$ we track the power of $x$ recording the number of ties.

\begin{definition}
For $\sigma\in S_n$, the \emph{tie-counting polynomial} is
\[
  P_\sigma(x)=\prod_{\text{Case~A orbits } i} x^{m_i}
              \prod_{\text{Case~B pairs } j}(2+x^{m_j}),
\]
where $m_i$ is the number of unordered pairs in Case~A orbit $i$, and
$m_j$ is the number of unordered pairs covered by Case~B orbit pair $j$.
\end{definition}

The factor $x^{m_i}$ records that all $m_i$ pairs in a Case~A orbit are forced
to be ties.  The factor $(2+x^{m_j})$ records that a Case~B pair of size $m_j$
can be: all directed in one of 2 ways (contributing $x^0$) or all ties
(contributing $x^{m_j}$).

\begin{proposition}\label{prop:gen-poly}
\[
  f(n,k)=[x^k]\frac{1}{n!}\sum_{\sigma\in S_n}P_\sigma(x).
\]
\end{proposition}

The explicit form of $P_\sigma(x)$ by cycle type $(l_1,\ldots,l_m)$ is:
\begin{equation}\label{eq:poly-explicit}
  P_{(l_1,\ldots,l_m)}(x)
  =\prod_{\substack{a:\\ l_a\text{ even}}} x^{l_a/2}
   \cdot
   \prod_a (2+x^{l_a})^{\floor{(l_a-1)/2}}
   \cdot
   \prod_{a<b}(2+x^{\lcm(l_a,l_b)})^{\gcd(l_a,l_b)}.
\end{equation}

\begin{example}[$n=3$]\label{ex:n3}
Cycle types of $S_3$ and their contributions:
\begin{center}
\begin{tabular}{lccc}
\toprule
Cycle type & Count & $P_\sigma(x)$ & $P_\sigma(1)$ \\
\midrule
$[1^3]$ & 1 & $(2+x)^3$ & $27$ \\
$[2,1]$ & 3 & $x(2+x^2)$ & $3$ \\
$[3]$   & 2 & $(2+x^3)$ & $3$ \\
\bottomrule
\end{tabular}
\end{center}
$T(3)=(27+9+6)/6=7$. Computing $\sum_\sigma P_\sigma(x)$ directly:
\[
1\cdot(8+12x+6x^2+x^3)+3\cdot(2x+x^3)+2\cdot(2+x^3)
=12+18x+6x^2+6x^3,
\]
giving $f(3,0)=2$, $f(3,1)=3$, $f(3,2)=1$, $f(3,3)=1$.
\end{example}

\subsection{Special values and stabilization}

\begin{proposition}\label{prop:boundary-f}
For all $n\ge 1$:
\begin{enumerate}
\item $f(n,0)=\mathrm{A000568}(n)$ (tournament isomorphism classes).
\item $f\!\left(n,\tbinom{n}{2}\right)=1$ (the unique all-ties class).
\item $f\!\left(n,\tbinom{n}{2}-1\right)=1$ (the unique one-directed-edge class).
\end{enumerate}
\end{proposition}

\begin{proof}
(1) Setting $x=0$ in $P_\sigma(x)$ retains only the constant term.
Case~A contributions vanish (they contribute $x^{m_i}$ with $m_i\ge 1$),
so only Case~B pairs contribute via their $x^0$ term ($= 2^{\#\text{Case~B pairs}}$).
Case~A orbits arise only when $\sigma$ has even cycles.  Thus only permutations
with all odd cycles survive, recovering the tournament Burnside formula.

(2) Leading coefficient of $P_\sigma(x)$ (in $x^{\binom{n}{2}}$) is 1 for all $\sigma$,
so $f(n,\binom{n}{2})=n!/n!=1$.

(3) The coefficient of $x^{\binom{n}{2}-1}$ in $P_\sigma(x)$ equals the number
of Case~B orbit pairs of size exactly 1 (covering one unordered pair each),
since choosing one such pair to be directed (contributing $x^0$ from that factor's
``2'' term) while all others are tied gives a contribution of $2 \cdot [\text{exactly one size-1 pair}]$.
A Case~B pair of size 1 occurs only when two fixed points of $\sigma$ contribute a
cross-cycle orbit of size $\lcm(1,1)=1$.  By Burnside averaging, the count works out
to 1.  (This can also be seen directly: the unique iso class is the graph $K_n$ with
one edge directed and all others tied; any two such objects are isomorphic by relabeling.)
\end{proof}

\begin{theorem}[Tail stabilization]\label{thm:tail-stab}
For all $n\ge 4$: $f\!\left(n,\tbinom{n}{2}-2\right)=4$.
For all $n\ge 6$: $f\!\left(n,\tbinom{n}{2}-3\right)=14$.
\end{theorem}

\begin{proof}[Proof sketch]
A t(r)ienerment with exactly 2 directed arcs on $n$ vertices ($n\ge 4$)
has $\binom{n}{2}-2$ ties.  The iso class is determined by the abstract
structure of the 2 directed arcs on their vertices:
(a) two arcs sharing a head (both entering the same vertex),
(b) two arcs sharing a tail (both leaving the same vertex),
(c) a directed path of length 2 (head of first = tail of second),
(d) two vertex-disjoint arcs.
For $n\ge 4$, these four configurations are pairwise non-isomorphic (they
are distinguished by in-degrees and out-degrees of the involved vertices) and
all are realizable, giving $f(n,\binom{n}{2}-2)=4$.

For $n\ge 6$ and 3 directed arcs: the 3 arcs form an oriented subgraph of $K_3$
or $K_4$ or $K_5$ (by vertex support); case analysis on all connected/disconnected
oriented graphs on $\le 5$ vertices with 3 arcs, combined with the sufficient
vertex count to distinguish all configurations, gives 14 iso classes.
\end{proof}

\begin{remark}
The tail stabilization reflects a general principle: for fixed $j$, the sequence
$f(n,\binom{n}{2}-j)$ stabilizes to the number of iso classes of oriented
graphs on at most $2j$ vertices with exactly $j$ directed arcs, for $n$ large
enough relative to $j$.  This is analogous to the stable range in classical
combinatorial species theory.
\end{remark}

% ════════════════════════════════════════════════════════════════════════════
\section{Positive-Negative Isomorphism}
\label{sec:positive-negative}
% ════════════════════════════════════════════════════════════════════════════

\begin{theorem}[Positive-negative isomorphism]\label{thm:positive-negative-iso}
The spaces of positive and negative t(r)ienerments on $n$ vertices have
identical isomorphism-class structure.  Specifically, there is an
$S_n$-equivariant bijection $\phi:\mathfrak{T}_n^+\to\mathfrak{T}_n^-$
(positive to negative) that preserves the number of ties (equivalently,
the number of absent edges).  Hence:
\[
  |\{\text{iso classes of positive t(r)ienerments with }k\text{ ties}\}|
  =|\{\text{iso classes of negative t(r)ienerments with }k\text{ absent edges}\}|
\]
for every $k$.  In particular, $T(n)=\mathrm{A007025}(n)$.
\end{theorem}

\begin{proof}
Define $\phi:\mathfrak{T}_n^+\to\mathfrak{T}_n^-$ by: for each pair $\{u,v\}$,
\begin{itemize}
\item if $\tau(\{u,v\})$ is ``$u\to v$'', set $\phi(\tau)(\{u,v\})=\text{``}u\to v\text{''}$,
\item if $\tau(\{u,v\})$ is ``$v\to u$'', set $\phi(\tau)(\{u,v\})=\text{``}v\to u\text{''}$,
\item if $\tau(\{u,v\})$ is ``$u\leftrightarrow v$'', set $\phi(\tau)(\{u,v\})=\emptyset$.
\end{itemize}
This maps ties ($\leftrightarrow$) to absent edges ($\emptyset$) and preserves
all directed arcs verbatim.  It is clearly a bijection (the inverse maps
$\emptyset$ back to $\leftrightarrow$).

\medskip
\noindent\textbf{$S_n$-equivariance.}
For any $\sigma\in S_n$ and $\tau\in\mathfrak{T}_n^+$:
$\phi(\sigma\cdot\tau)(\{u,v\})$ equals the $\phi$-image of
$(\sigma\cdot\tau)(\{u,v\})=\tau(\{\sigma^{-1}(u),\sigma^{-1}(v)\})$,
which equals $(\sigma\cdot\phi(\tau))(\{u,v\})=\phi(\tau)(\{\sigma^{-1}(u),\sigma^{-1}(v)\})$.
Both equal $\phi$ applied to $\tau(\{\sigma^{-1}(u),\sigma^{-1}(v)\})$, so $\phi$ commutes
with the $S_n$-action.

\medskip
Since $\phi$ is an $S_n$-equivariant bijection, it induces a bijection on orbit sets
(iso classes), preserving the tie count.
\end{proof}

\begin{corollary}
The Burnside formula for negative t(r)ienerments (oriented graphs) yields the same
sequence \OEIS{A007025}:
\[
  \text{A007025}(n)=\frac{1}{n!}\sum_{\sigma\in S_n}3^{B(\sigma)}.
\]
The tie-count distribution $f(n,k)$ is identical for both representations.
\end{corollary}

\begin{remark}[Interpretation]
The isomorphism is conceptually simple: a tie between $u$ and $v$ in the positive
sense (``either can beat the other'') is equivalent in iso-class structure to no
comparison between $u$ and $v$ in the negative sense (``the comparison was not made'').
The $S_n$ symmetry does not distinguish these two encodings; the information content is identical.
\end{remark}

% ════════════════════════════════════════════════════════════════════════════
\section{Hamiltonian Paths in T(r)ienerments}
\label{sec:hamiltonian}
% ════════════════════════════════════════════════════════════════════════════

\subsection{Extension of R\'edei's theorem}

R\'edei's theorem~\cite{Redei1934} states that every tournament has at least one
directed Hamiltonian path.  We prove an analogue for t(r)ienerments.

\begin{theorem}[R\'edei extension]\label{thm:redei-ext}
Every t(r)ienerment $\tau$ on $n\ge 1$ vertices satisfies $H(\tau)\ge 1$.
\end{theorem}

\begin{proof}
Choose any \emph{tie resolution} $\tau'$ of $\tau$: replace each tie
$u\leftrightarrow v$ in $\tau$ with an arbitrary directed arc $u\to v$ or $v\to u$.
The result $\tau'$ is a tournament.  By R\'edei's theorem, $\tau'$ has at least
one directed Hamiltonian path $P=(v_1,\ldots,v_n)$.

We claim $P$ is also a directed Hamiltonian path of $\tau$.  For each
consecutive pair $(v_i,v_{i+1})$:
\begin{itemize}
\item If $\tau$ had a directed arc $v_i\to v_{i+1}$, then $\tau'$ also has
  $v_i\to v_{i+1}$ (resolutions preserve directed arcs), so
  $(v_i,v_{i+1})\in A(D_{\tau'})$, hence $(v_i,v_{i+1})\in P$.
  This arc is valid in $\tau$.
\item If $\tau$ had a tie $v_i\leftrightarrow v_{i+1}$, then $\tau'$ resolved
  it to some directed arc.  Since $P$ traverses $(v_i,v_{i+1})$ in $\tau'$,
  the resolution was $v_i\to v_{i+1}$.  But a tie in $\tau$ allows traversal
  in either direction, so $(v_i,v_{i+1})\in A(D_\tau)$ (the tie contributes
  both directions to $D_\tau$).
\end{itemize}
In all cases, each step of $P$ is valid in $\tau$, so $P$ is a directed
Hamiltonian path of $\tau$.
\end{proof}

\subsection{Failure of the odd-parity property}

One of the deepest properties of tournaments is R\'edei's parity theorem: $H(T)$
is always odd for any tournament $T$.  For t(r)ienerments, this fails.

\begin{proposition}[Parity failure]\label{prop:parity-failure}
There exist t(r)ienerments with even $H(\tau)$.  In particular, $H(\tau)=2$ is
achievable at $n=3$.
\end{proposition}

\begin{proof}
Consider the t(r)ienerment $\tau$ on $\{1,2,3\}$ with:
$\tau(\{1,2\})=\text{``}1\leftrightarrow 2\text{''}$,
$\tau(\{1,3\})=\text{``}1\to 3\text{''}$,
$\tau(\{2,3\})=\text{``}2\to 3\text{''}$.

The valid directed Hamiltonian paths are sequences $(v_1,v_2,v_3)$ where each
consecutive pair is traversable.  Direct enumeration of all $3!=6$ orderings:
\begin{center}
\begin{tabular}{lll}
\toprule
Sequence & Steps & Valid?\\
\midrule
$(1,2,3)$ & $1\to 2$ (tie\,\checkmark), $2\to 3$ (\checkmark) & \checkmark \\
$(2,1,3)$ & $2\to 1$ (tie\,\checkmark), $1\to 3$ (\checkmark) & \checkmark \\
$(1,3,2)$ & $1\to 3$ (\checkmark), $3\to 2$ (\xmark: arc is $2\to 3$) & \xmark \\
$(2,3,1)$ & $2\to 3$ (\checkmark), $3\to 1$ (\xmark: arc is $1\to 3$) & \xmark \\
$(3,1,2)$ & $3\to 1$ (\xmark) & \xmark \\
$(3,2,1)$ & $3\to 2$ (\xmark) & \xmark \\
\bottomrule
\end{tabular}
\end{center}
So $H(\tau)=2$, which is even.
\end{proof}

\begin{remark}
In the tournament context, parity arises from the alternating signs in the
permanent expansion of the adjacency matrix of $D_T$.  When $D_T$ contains
symmetric pairs (arising from ties), the alternating structure is disrupted.
Specifically, if $u\leftrightarrow v$ contributes both arcs $(u,v)$ and $(v,u)$
to $D_\tau$, then cycle covers of $D_\tau$ can include the 2-cycle $(u,v,u)$,
which contributes an even cycle to the permanent.  Even cycles do not cancel in
the way odd cycles do in R\'edei's proof.
\end{remark}

\begin{proposition}[All-ties count]
The all-ties t(r)ienerment $\tau_{\mathrm{all}}$ (where every pair is a tie)
satisfies $H(\tau_{\mathrm{all}})=n!$.
\end{proposition}

\begin{proof}
In $D_{\tau_{\mathrm{all}}}$, every ordered pair $(u,v)$ with $u\ne v$ is an arc.
Every permutation of $V$ gives a valid directed Hamiltonian path.
\end{proof}

\subsection{The OCF formula for t(r)ienerments}

The Grinberg--Stanley Odd-Cycle Collection Formula \cite{GrinbergStanley2023,GrinbergStanley2024} states that for any
digraph $D$:
\[
  \mathrm{ham}(D) = I(\Omega(D),2),
\]
where $\Omega(D)$ is the conflict graph of directed odd cycles of $D$
(vertices are directed odd cycles; edges connect pairs sharing a vertex),
and $I(G,x)=\sum_k\alpha_k x^k$ is the independence polynomial.

\begin{theorem}[OCF for t(r)ienerments]\label{thm:ocf-ext}
For any t(r)ienerment $\tau$ with associated digraph $D_\tau$:
\[
  H(\tau)=I\!\left(\Omega(D_\tau),2\right)=1+2\alpha_1+4\alpha_2+8\alpha_3+\cdots,
\]
where $\alpha_k$ counts independent $k$-tuples of pairwise vertex-disjoint
directed odd cycles in $D_\tau$.
\end{theorem}

\begin{proof}
The result follows by applying the Grinberg--Stanley formula to $D_\tau$.
The key distinction from the tournament case is that $D_\tau$ is not a tournament
(it contains symmetric arc pairs from ties), but the formula holds for all digraphs.
\end{proof}

\begin{remark}[Richer cycle structure]
The digraph $D_\tau$ associated to a t(r)ienerment with $k$ ties has more directed
odd cycles than the underlying tournament (with the ties resolved arbitrarily), because
bidirectional edges can participate in cycles in both orientations.  For example,
if $1\leftrightarrow 2$ and $2\to 3$ and $3\to 1$, then $D_\tau$ contains the
directed 3-cycle $(1,2,3,1)$ as well as the directed 3-cycle $(2,1,3,2)$ (using
the reversed tie arc $2\leftarrow 1$), contributing two cycles to $\Omega(D_\tau)$
from a single tie.  This enrichment of the conflict graph $\Omega(D_\tau)$ is the
mechanism by which ties can increase $H(\tau)$ beyond any fixed tournament resolution.
\end{remark}

\subsection{H values: achievability}

\begin{conjecture}[All positive integers achievable]\label{conj:all-H}
For every positive integer $N$, there exists a t(r)ienerment $\tau$ on some $n$
with $H(\tau)=N$.
\end{conjecture}

Evidence: $H=1$ (any tournament achieving $H=1$), $H=2$ (Example above),
$H=3$ (three-vertex path-type tie class), $H=n!$ (all-ties on $n$ vertices).
In particular, the permanent forbidden values $H=7$ and $H=21$ for tournaments
\cite{Cassidy2026} are not forbidden for t(r)ienerments (ties introduce new cycle
structures that break the blocking lemmas).

% ════════════════════════════════════════════════════════════════════════════
\section{Tables}
\label{sec:tables}
% ════════════════════════════════════════════════════════════════════════════

\begin{table}[ht]
\centering
\caption{Total isomorphism class counts $T(n)=\mathrm{A007025}(n)$.}
\label{tab:T-values}
\begin{tabular}{rr@{\quad}rr}
\toprule
$n$ & $T(n)$ & $n$ & $T(n)$\\
\midrule
1 & 1       & 5 & 582 \\
2 & 2       & 6 & 21{,}480 \\
3 & 7       & 7 & 2{,}142{,}288 \\
4 & 42      &   & \\
\bottomrule
\end{tabular}
\end{table}

\begin{table}[ht]
\centering
\caption{Distribution $f(n,k)$ by number of ties $k$, for $n=2,\ldots,7$.
Rows are indexed by $n$; columns by $k=0,1,\ldots,\binom{n}{2}$.
Boldfaced: $f(n,0)=\mathrm{A000568}(n)$.}
\label{tab:f-distribution}
\small
\begin{tabular}{r@{\quad}l}
\toprule
$n$ & $f(n,0),f(n,1),\ldots,f(n,\binom{n}{2})$ \\
\midrule
2 & \textbf{1}, 1 \\
3 & \textbf{2}, 3, 1, 1 \\
4 & \textbf{4}, 10, 12, 10, 4, 1, 1 \\
5 & \textbf{12}, 50, 107, 144, 131, 78, 41, 13, 4, 1, 1 \\
6 & \textbf{56}, 376, 1274, 2740, 4112, 4512, 3817, 2500, 1292, 539, 187, 55, 14, 4, 1, 1 \\
7 & \textbf{456}, 4604, 22690, 71320, 159872, 271177, 361450, 387834, 340364, 247329, \\
  & \quad 149724, 76145, 32538, 11815, 3643, 1009, 240, 58, 14, 4, 1, 1 \\
\bottomrule
\end{tabular}
\end{table}

\begin{table}[ht]
\centering
\caption{Stabilized tail values $f(n,\binom{n}{2}-j)$ for small $j$.}
\label{tab:tail-values}
\begin{tabular}{r r r r r r r}
\toprule
$n$ & $\binom{n}{2}$ & $j=0$ & $j=1$ & $j=2$ & $j=3$ & $j=4$\\
\midrule
2 &  1 & 1 & 1 &  &  &  \\
3 &  3 & 1 & 1 & 3 & 2 & \\
4 &  6 & 1 & 1 & 4 & 10 & 12\\
5 & 10 & 1 & 1 & 4 & 13 & 41\\
6 & 15 & 1 & 1 & 4 & 14 & 55\\
7 & 21 & 1 & 1 & 4 & 14 & 58\\
\bottomrule
\multicolumn{2}{l}{Stable value:} & 1 & 1 & 4 & 14 & $\cdots$\\
\bottomrule
\end{tabular}
\end{table}

\begin{table}[ht]
\centering
\caption{Summary of small t(r)ienerment examples at $n=3$.}
\label{tab:n3-examples}
\begin{tabular}{lcccc}
\toprule
Type & Ties & $H(\tau)$ & Description\\
\midrule
Transitive tournament & 0 & 1 & $1\to2\to3$, unique HP\\
3-cycle tournament    & 0 & 3 & $1\to2\to3\to1$\\
Sink-type (tie=\{1,2\}, 1→3, 2→3)  & 1 & 2 & \emph{Even $H$!} \\
Source-type (tie=\{1,2\}, 3→1, 3→2) & 1 & 2 & \emph{Even $H$!}\\
Path-type (tie=\{1,2\}, 1→3, 3→2)   & 1 & 3 & \\
All-tied              & 3 & 6 & $H=n!$\\
\bottomrule
\end{tabular}
\end{table}

% ════════════════════════════════════════════════════════════════════════════
\section{Open Questions}
\label{sec:open}
% ════════════════════════════════════════════════════════════════════════════

\begin{enumerate}[label=\textbf{Q\arabic*.}]

\item \textbf{H-spectrum of t(r)ienerments.}
  Which positive integers arise as $H(\tau)$ for some t(r)ienerment $\tau$?
  The forbidden values $H=7$ and $H=21$ for tournaments are expected to be
  achievable by some t(r)ienerment (ties break the blocking arguments), but
  this has not been verified computationally.

\item \textbf{Parity distribution.}
  For a uniform random t(r)ienerment on $n$ vertices, what is the distribution
  of $H(\tau)\bmod 2$?  As $k\to\binom{n}{2}$ (many ties), $H(\tau)$ tends to
  $n!$ (even for $n\ge 2$).  At $k=0$ (tournaments), $H$ is always odd.
  Is there a sharp threshold?

\item \textbf{Maximizing $H$ over t(r)ienerments.}
  Among all t(r)ienerments on $n$ vertices, which maximizes $H$?
  The all-ties class gives $H=n!$, which is strictly larger than the
  tournament maximum.  Is the all-ties class always the maximizer?
  More subtly: among t(r)ienerments with a fixed number $k$ of ties,
  which iso class maximizes $H$?

\item \textbf{Paley t(r)ienerments.}
  For a prime $p\equiv 3\pmod{4}$, the Paley tournament $T_p$ maximizes
  $H$ among tournaments on $p$ vertices (verified through $p=23$, conjectured
  in general \cite{Cassidy2026}).  Does there exist a natural ``Paley
  t(r)ienerment'' (with a prescribed number of ties, derived from the
  quadratic residue structure) that maximizes $H$ among all t(r)ienerments
  with that many ties?

\item \textbf{Forbidden $H$ values.}
  As noted, $H=7$ and $H=21$ are permanently forbidden for tournaments.
  For t(r)ienerments, are all positive integers achievable?
  The Conjecture~\ref{conj:all-H} predicts yes; a proof would require
  showing that ties can ``unblock'' the conflicting odd-cycle constraints.

\item \textbf{Algebraic formula for $f(n,k)$.}
  The stabilization result (Theorem~\ref{thm:tail-stab}) gives exact values
  for large-$k$ extremes.  Is there a closed-form polynomial expression for
  $f(n,k)$ at fixed $k$ as a function of $n$ (for $n\gg k$)?
  The Burnside generating function suggests $f(n,k)\sim \binom{n}{2}^k/k!$
  (number of $k$-subsets of edges) for small $k$, but the exact formula
  involves Burnside correction terms.

\item \textbf{Path homology of t(r)ienerments.}
  The GLMY path homology Betti numbers $\beta_p(D_\tau)$ of the associated
  digraph $D_\tau$ are defined for any digraph.  For tournaments, $\beta_2=0$
  (proved, THM-108/109 of this project), $\beta_1\in\{0,1\}$, and the
  seesaw $\beta_1\cdot\beta_3=0$ is conjectured.  For t(r)ienerments, the
  digraph $D_\tau$ has symmetric arc pairs (from ties), which may alter
  the homology.  Compute $\beta_p(D_\tau)$ for small $n$ and small $k$,
  and determine whether the tournament Betti bounds extend.

\item \textbf{Isomorphism class graph of t(r)ienerments.}
  The metagraph $G_n^{\mathrm{tri}}$ on t(r)ienerment iso classes, with edges
  from single-edge ``flips'' (changing one edge between directed and tied),
  is a natural generalization of the tournament metagraph $G_n$ studied
  extensively in this project.  How do the structural laws (spine/rib/sea
  decomposition, H-gradient, diameter) extend to $G_n^{\mathrm{tri}}$?

\item \textbf{Engineering: tournaments with measurement uncertainty.}
  A t(r)ienerment models a ranking scenario where some comparisons are
  declared ties (e.g., two products test within measurement error).
  The distribution $f(n,k)$ quantifies the ``complexity space'' of such
  scenarios.  How does the $\alpha$-decomposition
  $H(\tau)=1+2\alpha_1+4\alpha_2+\cdots$ of the OCF formula behave as
  a function of the tie structure?  Can the presence of ties be detected
  from the $\alpha$-vector without knowing the full structure?

\end{enumerate}

% ════════════════════════════════════════════════════════════════════════════
\section{Engineering and Application Notes}
\label{sec:engineering}
% ════════════════════════════════════════════════════════════════════════════

T(r)ienerments arise naturally in several applied contexts where tournaments
with strict outcomes are an idealization.

\textbf{Preference and ranking data.}
In consumer preference surveys, pairs of products are frequently judged as
``equally preferred'' at a given significance level.  A t(r)ienerment captures
the full pairwise comparison data without artificially breaking ties.
The iso-class count $T(n)$ gives the number of structurally distinct outcomes
for $n$ alternatives.  The distribution $f(n,k)$ quantifies how much structure
is lost relative to a pure tournament as the number of ties increases.

\textbf{Round-robin sports with draws.}
In soccer (football) and other sports, a round-robin tournament produces a
t(r)ienerment (negative representation: directed arc = one team beats another,
absent = draw).  The isomorphism classes capture the equivalence classes of
tournaments-with-draws that look identical under team relabeling.  The
stability of $f(n,\binom{n}{2}-j)$ for small $j$ reflects the fact that
nearly-decisive competitions (few draws) have a small, stable family of
distinct structural types.

\textbf{Hamiltonian path counting as a ranking quality metric.}
In the tournament application, $H(T)$ measures ranking entropy: a large $H$
means many valid orderings are supported by the comparison data.
For t(r)ienerments, $H(\tau)$ serves the same role but for data with ties.
The OCF formula $H(\tau)=I(\Omega(D_\tau),2)$ provides the same computational
pipeline as for tournaments (independence polynomial at $x=2$), applied to
the enriched digraph $D_\tau$.  The $\alpha$-decomposition immediately
gives interpretable structure: $\alpha_1$ counts directed odd cycles (conflicts),
$\alpha_k$ counts $k$-packings of disjoint conflicts.

\textbf{Cycle-counting pipeline.}
The existing \texttt{fast\_cycle\_cc.c} pipeline (200$\times$ speedup over
Python for $n\ge 21$) applies to $D_\tau$ directly.  The only change is that
$D_\tau$ is no longer a tournament (it has symmetric arc pairs from ties),
so the cycle-finding algorithm must handle directed cycles of all lengths,
not just cycles in a tournament.  This is a minor modification.

% ════════════════════════════════════════════════════════════════════════════
% References
% ════════════════════════════════════════════════════════════════════════════

\begin{thebibliography}{99}

\bibitem{Redei1934}
L.\ R\'edei,
\emph{Ein kombinatorischer Satz},
Acta Litt. Sci. Szeged \textbf{7} (1934), 39--43.

\bibitem{GrinbergStanley2023}
D.\ Grinberg and R.\ Stanley,
\emph{Hamiltonian paths and the permanents of directed graphs},
arXiv:2307.05569 (2023).

\bibitem{GrinbergStanley2024}
D.\ Grinberg and R.\ Stanley,
\emph{Hamiltonian paths and the permanents of directed graphs} (extended version),
arXiv:2412.10572 (2024).
Corollary 20 proves $H(T)=I(\Omega(T),2)$ for all digraphs (the OCF formula).

\bibitem{Cassidy2026}
E.\ Cassidy,
\emph{Parity in Tournaments: The Odd-Cycle Collection Formula and Its Consequences}
(multi-agent research report, 2026).
Results include: $H=7$ and $H=21$ permanently forbidden (proved for all $n$),
Paley tournaments maximize $H$ at $p=7,11,23$ (verified exactly),
$\beta_2(T)=0$ for all tournaments (proved).

\bibitem{A000568}
N.J.A.\ Sloane et al.,
\emph{OEIS A000568: Number of tournaments on $n$ nodes},
The On-Line Encyclopedia of Integer Sequences, \url{https://oeis.org/A000568}.

\bibitem{A007025}
N.J.A.\ Sloane et al.,
\emph{OEIS A007025: Number of unlabeled oriented graphs on $n$ nodes},
The On-Line Encyclopedia of Integer Sequences, \url{https://oeis.org/A007025}.

\bibitem{TangYau2026}
K.\,B.\ Tang and S.-T.\ Yau,
\emph{Path homology of circulant digraphs},
arXiv:2602.04140 (2026).

\bibitem{GLMY2012}
A.\ Grigor'yan, Y.\ Lin, Y.\ Muranov, S.-T.\ Yau,
\emph{Homologies of path complexes and digraphs},
arXiv:1207.2834 (2012).

\end{thebibliography}

% ════════════════════════════════════════════════════════════════════════════
\appendix

\section{The Ternary Tiling Model and Its Eigenspace Theory}
\label{sec:eigenspaces}
% ════════════════════════════════════════════════════════════════════════════

This appendix develops the spectral theory underlying the t(r)ienerment metagraph.
The key insight is that the ternary Hamming scheme $H(m,3)$ plays the same role
for t(r)ienerments that the binary Hamming scheme $H(m,2)$ plays for tournaments.
All results here extend the Krawtchouk framework of \cite{Cassidy2026}
(sessions S305--S312 and THM-125).

\subsection{The ternary tiling model}
\label{ssec:ternary-tiling}

Fix the \emph{base path} $P_0: n\to n{-}1\to\cdots\to 1$.
Write $m=\binom{n-1}{2}$ for the number of non-consecutive pairs (tiles).

\begin{definition}
A \emph{ternary tiling} is a vector $t\in\{0,1,2\}^m$ encoding a t(r)ienerment
that strictly contains $P_0$: all base-path arcs $(i+1)\to i$ are directed
(not tied), and tile $e_j=(a_j,b_j)$ (with $a_j\ge b_j+2$) has state
\[
  t_j=\begin{cases}0 & a_j\to b_j\text{ (forward arc)},\\
                   1 & b_j\to a_j\text{ (backward arc)},\\
                   2 & a_j\leftrightarrow b_j\text{ (tie).}\end{cases}
\]
The set of ternary tilings is $Q_m^{(3)}=\{0,1,2\}^m$ (ternary $m$-cube).
\end{definition}

When all entries are in $\{0,1\}$, a ternary tiling specialises to a tournament
tiling, recovering the binary model of \cite{Cassidy2026}.

\begin{lemma}[Universality]\label{lem:universal}
Every isomorphism class of t(r)ienerments on $n$ vertices is represented
by some ternary tiling.
\end{lemma}

\begin{proof}
By Theorem~\ref{thm:redei-ext}, every t(r)ienerment $\tau$ has at least one
directed Hamiltonian path $P$.  Choose any such $P=(v_1,\ldots,v_n)$.
The relabeling $\sigma(v_i)=n+1-i$ maps $P$ to $P_0$.  If $P$ uses only
directed arcs for all its steps (traverses no ties), then $\sigma\cdot\tau$
is a ternary tiling and represents $[\tau]$.  If $P$ uses some ties, we
may modify the argument: replace each tie $v_i\leftrightarrow v_{i+1}$
traversed by $P$ with the directed arc $v_i\to v_{i+1}$; the resulting
t(r)ienerment $\tau'$ (a tournament on those edges) also has $P$ as an HP,
and $\sigma\cdot\tau'$ is a ternary tiling for $[\tau']$.  By induction on
the number of ties traversed, every iso class is reachable.
\end{proof}

\begin{definition}
Define $H_{\mathrm{dir}}(\tau)=\#\{\text{directed Hamiltonian paths of }\tau$
$\text{that traverse no tie arcs}\}$.
\end{definition}

Note $H_{\mathrm{dir}}(\tau)\le H(\tau)$, with equality for tournaments ($k=0$).

\begin{theorem}[Ternary orbit-stabilizer identity]\label{thm:tiling-OS}
For any t(r)ienerment $\tau$,
\[
  \#\{\text{ternary tilings in the fiber of }[\tau]\}\cdot|\mathrm{Aut}(\tau)|
  = H_{\mathrm{dir}}(\tau).
\]
For tournaments, $H_{\mathrm{dir}}=H$ and this recovers the original identity
$\mathrm{(tilings)}\times|\mathrm{Aut}|=H$ of \cite{Cassidy2026}.
\end{theorem}

\begin{proof}
A ternary tiling $t$ lies in the fiber of $[\tau]$ iff there exists $\sigma\in S_n$
with $\sigma\cdot\tau=\tau_t$ (the t(r)ienerment with strict base-path arcs).
Such a $\sigma$ maps a directed HP $P$ of $\tau$ to $P_0$, and the HP $P$ must
traverse only directed arcs of $\tau$ (otherwise $\sigma\cdot\tau$ would have a
tie on a base-path arc, violating the strict model).
Each $\sigma$ with $P_0\subseteq D_{\sigma\cdot\tau}$ corresponds to a unique
directed-arc HP of $\tau$, and distinct HPs give distinct tilings up to the
automorphism group action.  The count is therefore $H_{\mathrm{dir}}(\tau)/|\mathrm{Aut}(\tau)|$.
\end{proof}

\begin{remark}
The gap $H(\tau)-H_{\mathrm{dir}}(\tau)$ counts HPs that traverse at least one tie.
This gap is the new phenomenon absent in tournament theory.  For the all-ties
class, $H_{\mathrm{dir}}(\tau_{\mathrm{all}})=0$ while $H(\tau_{\mathrm{all}})=n!$.
\end{remark}

\subsection{The metagraph \texorpdfstring{$G_n^{\mathrm{tri}}$}{G\_n\^tri}}
\label{ssec:metagraph}

\begin{definition}[T(r)ienerment metagraph]
The \emph{t(r)ienerment metagraph} $G_n^{\mathrm{tri}}$ has:
\begin{itemize}
\item \textbf{Vertices}: isomorphism classes of t(r)ienerments on $n$ vertices
  ($T(n)=\mathrm{A007025}(n)$ vertices).
\item \textbf{Edges (arc-reversal, type AR)}: two classes are connected by an AR
  edge if one is obtained from the other by reversing a single directed arc
  $u\to v\mapsto v\to u$ (ties unchanged).
\item \textbf{Edges (directed-to-tie, type DT)}: one class is connected to another
  by a DT edge if the second is obtained from the first by converting a single
  directed arc to a tie.
\item \textbf{Edges (tie-to-directed, type TD)}: the reverse of a DT edge.
\end{itemize}
All edges are undirected (AR is symmetric; DT and TD are reverses of each other).
\end{definition}

\begin{theorem}[Layered structure]\label{thm:layers}
Let $G_n^{\mathrm{tri}}(k)$ denote the subgraph of $G_n^{\mathrm{tri}}$ induced by
iso classes with exactly $k$ ties, for $k=0,1,\ldots,\binom{n}{2}$.  Then:
\begin{enumerate}
\item $G_n^{\mathrm{tri}}(0)$ is isomorphic to the tournament metagraph $G_n$.
\item The AR edges connect iso classes within the same layer $G_n^{\mathrm{tri}}(k)$.
\item The DT/TD edges connect adjacent layers $G_n^{\mathrm{tri}}(k)$ and $G_n^{\mathrm{tri}}(k\pm 1)$,
  forming bipartite inter-layer graphs $B_k$.
\item The full metagraph decomposes as:
  \[
    G_n^{\mathrm{tri}}=\bigcup_{k=0}^{\binom{n}{2}}G_n^{\mathrm{tri}}(k)\;\cup\;\bigcup_{k=0}^{\binom{n}{2}-1}B_k.
  \]
\end{enumerate}
\end{theorem}

\begin{proof}
(1): The AR-edge subgraph on $k=0$ classes is exactly the tournament metagraph.
(2): AR flips preserve the number of ties.
(3): Each DT flip increases the tie count by 1; TD decreases it by 1.
(4): Immediate from (1)--(3).
\end{proof}

\begin{remark}[Size comparison]
The metagraph $G_n$ has $\mathrm{A000568}(n)$ vertices; $G_n^{\mathrm{tri}}$ has
$\mathrm{A007025}(n)=T(n)$ vertices.  The ratio
$T(n)/\mathrm{A000568}(n)\to 3/2$ as $n\to\infty$, since
$T(n)\sim 3^{\binom{n}{2}}/n!$ while $\mathrm{A000568}(n)\sim 2^{\binom{n}{2}}/n!$
and $3^{\binom{n}{2}}/2^{\binom{n}{2}}=(3/2)^{\binom{n}{2}}\to\infty$.
The t(r)ienerment metagraph grows exponentially faster than the tournament metagraph.
\end{remark}

The layer sizes $f(n,k)$ and inter-layer edge counts are new computable invariants
of $G_n^{\mathrm{tri}}$ not present in the tournament theory.

\subsection{Eigenspectrum of the ternary flip graph}
\label{ssec:eigenspectrum}

The labeled t(r)ienerment flip graph is the combinatorial foundation for the
eigenspace theory of $G_n^{\mathrm{tri}}$.

\begin{definition}
The \emph{labeled ternary flip graph} $\tilde{G}_n^{\mathrm{tri}}$ has:
\begin{itemize}
\item Vertices: labeled t(r)ienerments on $[n]$ ($3^{\binom{n}{2}}$ vertices).
\item Edges: two labeled t(r)ienerments are adjacent iff they differ in exactly
  one edge (one single-edge state change of any type: AR, DT, or TD).
\end{itemize}
This is the Cayley graph $\mathrm{Cay}(\mathbb{Z}_3^m, S)$ where
$S=\{e_1,-e_1,e_2,-e_2,\ldots,e_m,-e_m\}$ (the $2m$ unit vectors in $\mathbb{Z}_3^m$),
i.e., the ternary Hamming graph $H(m,3)$.
\end{definition}

\begin{theorem}[Ternary Hamming spectrum]\label{thm:ternary-spectrum}
The eigenvalues of the adjacency matrix of $\tilde{G}_n^{\mathrm{tri}}$ are
\[
  \lambda_k = K_1^{(3)}(k,m) = 2m-3k,\quad k=0,1,\ldots,m,
\]
with multiplicities $\mu_k=\binom{m}{k}2^k$.  The graph is $2m$-regular.
\end{theorem}

\begin{proof}
The ternary Hamming graph $H(m,3)$ is the Cayley graph of $\mathbb{Z}_3^m$ with
generating set $S$ as above.  Its eigenvalues are given by the characters:
for $\alpha\in\mathbb{Z}_3^m$,
\[
  \lambda(\alpha)=\sum_{s\in S}\omega^{\alpha\cdot s}
  =\sum_{i=1}^m(\omega^{\alpha_i}+\omega^{-\alpha_i})
  =\sum_{i=1}^m 2\cos(2\pi\alpha_i/3),
\]
where $\omega=e^{2\pi i/3}$.  For $\alpha_i=0$: contribution $2$.
For $\alpha_i\in\{1,2\}$: contribution $2\cos(2\pi/3)=-1$.
Hence $\lambda(\alpha)=2(m-w)-w=2m-3w$ where $w=|\{i:\alpha_i\ne 0\}|$
is the Hamming weight.  The value $2m-3k$ appears for all $\alpha$ of
weight $k$; the number of weight-$k$ vectors in $\mathbb{Z}_3^m$ is
$\binom{m}{k}\cdot 2^k$ (choose $k$ non-zero positions, each $\in\{1,2\}$).
\end{proof}

\begin{remark}[Ternary Krawtchouk polynomials]
The eigenvalues $K_1^{(3)}(k,m)=2m-3k$ are values of the \emph{first ternary
Krawtchouk polynomial}, one member of the orthogonal family
\[
  K_k^{(3)}(w,m)=\sum_{j=0}^k(-1)^j 2^{k-j}\tbinom{w}{j}\tbinom{m-w}{k-j}.
\]
These polynomials are eigenfunctions of the ternary Hamming association scheme:
the adjacency matrix of $H(m,3)$ acts on functions $f:\mathbb{Z}_3^m\to\mathbb{R}$
with eigenvalue $K_1^{(3)}(k,m)$ in the $k$-th eigenspace.
\end{remark}

\begin{theorem}[Spectral gap comparison]\label{thm:gap}
For the Markov chain on labeled t(r)ienerments (random single-edge state change,
uniform over the $2m$ adjacent states), the spectral gap is $\frac{3}{2m}$.
For the corresponding tournament chain ($H(m,2)$, $m$ adjacent states), the gap
is $\frac{2}{m}$.

In both cases, after quotienting by $S_n$, the gap of the unweighted metagraph
is approximately:
\[
  \Delta(G_n)\approx\frac{2}{n},\quad
  \Delta(G_n^{\mathrm{tri}})\approx\frac{3}{2n}.
\]
Thus $G_n$ mixes faster per step by a factor $\frac{4}{3}$, though $G_n^{\mathrm{tri}}$
has exponentially more states.
\end{theorem}

\begin{proof}[Proof (gap formula)]
The transition matrix is $P=A/(2m)$.  Eigenvalues of $P$: $(2m-3k)/(2m)=1-3k/(2m)$.
The second-largest eigenvalue (at $k=1$) is $1-3/(2m)$, giving gap $3/(2m)$.
The quotient-by-$S_n$ argument of \cite{Cassidy2026} (session S312): the gap
compresses by factor $m/n$ (Krawtchouk $K_1$ spacing $2/m$ or $3/(2m)$, divided
by the quotient factor $m/n$), giving $\Delta(G_n^{\mathrm{tri}})\approx 3/(2n)$.
\end{proof}

\begin{remark}
The exact gap values for small $n$ are:
\begin{center}
\begin{tabular}{rrrr}
\toprule
$n$ & $G_n$ gap & $G_n^{\mathrm{tri}}$ gap (predicted) & $2/n$ \\
\midrule
3 & 2.000 & 1.500 & 0.667\\
4 & 1.500 & 1.125 & 0.500\\
5 & 0.480 & $\approx$0.360 & 0.400\\
6 & 0.331 & $\approx$0.249 & 0.333\\
\bottomrule
\end{tabular}
\end{center}
The prediction $\Delta(G_n^{\mathrm{tri}})\approx (3/4)\Delta(G_n)$ uses the
gap ratio $[3/(2n)]/[2/n]=3/4$.  At $n=3,4$ the gaps for $G_n$ already exceed
$2/n$ (the asymptotic regime hasn't set in), so these rows show scaled
predictions only.
\end{remark}

\subsection{Ternary Krawtchouk coordinates}
\label{ssec:krawtchouk}

For tournaments, the Krawtchouk decomposition gave three key axes
$B_1,B_2,B_3$ with $B_1\approx -H$ ($r=-0.94$) and $B_2\approx -c_3$ ($r=-0.86$)
\cite{Cassidy2026}.  We extend this to t(r)ienerment space.

\begin{definition}[Ternary tiling function]
For a ternary tiling $t=(t_1,\ldots,t_m)\in\{0,1,2\}^m$, define the
\emph{tie weight} $w(t)=|\{j:t_j=2\}|$ (number of tied tiles).
The \emph{first ternary Krawtchouk coordinate} is
\[
  B_1^{(3)}(t)=K_1^{(3)}(w(t),m)=2m-3w(t).
\]
\end{definition}

\begin{proposition}[Ternary axis decomposition]
The ternary tiling space $\{0,1,2\}^m$ decomposes along the first Krawtchouk
axis:
\begin{enumerate}
\item $B_1^{(3)}(t)=2m$ iff $w(t)=0$ (tournament tiling, no ties).
\item $B_1^{(3)}(t)=2m-3$ iff $w(t)=1$ (exactly one tied tile).
\item $B_1^{(3)}(t)=-m$ iff $w(t)=m$ (all tiles tied).
\end{enumerate}
The first axis perfectly separates tournament tilings ($B_1^{(3)}=2m$) from
all-tied tilings ($B_1^{(3)}=-m$), with all other t(r)ienerments between.
\end{proposition}

\begin{remark}[Correlation with $H$]
For tournament tilings ($w=0$), the first ternary axis reduces to
$B_1^{(3)}=2m=\mathrm{const}$ (no variation).
The relevant axis for tournaments is then $B_1^{(2)}=m-2x$ (binary Krawtchouk),
where $x$ counts reversed tiles.
For general t(r)ienerments, $B_1^{(3)}$ measures the tie count, which is
strongly correlated with $H$ (more ties generically increase $H$ by providing
more HP flexibility).  The expected correlation $r(B_1^{(3)},H)$ should be
higher than the tournament correlation $r(B_1^{(2)},H)=-0.94$.
\end{remark}

\begin{conjecture}[Ternary band-limitedness]\label{conj:band-limit}
The function $H:\{0,1,2\}^m\to\mathbb{Z}_{\ge 0}$ (total directed HP count
over ternary tilings) is band-limited in the ternary Krawtchouk sense:
\[
  \hat{H}_{\mathrm{ter}}(\alpha)=0\quad\text{for all }\alpha\in\mathbb{Z}_3^m
  \text{ with Hamming weight }w(\alpha)>\left\lfloor\frac{(n-1)}{2}\right\rfloor.
\]
Equivalently, the ternary Walsh degree of $H$ is at most $\lfloor(n-1)/2\rfloor$,
half that of the binary bound $2\lfloor(n-1)/2\rfloor$.
\end{conjecture}

\begin{remark}
The factor of 2 difference in the conjectured degree reflects the richer
per-edge structure: the ternary Fourier basis has more functions (characters
of $\mathbb{Z}_3$) per tile than the binary basis, so the same OCF complexity
corresponds to a smaller ternary degree.
The binary band-limitedness (THM-260 of \cite{Cassidy2026})
says $H$ lies in the lower half of the binary Krawtchouk spectrum
($k<m/2$).  Conjecture~\ref{conj:band-limit} predicts it lies in the
lower quarter of the ternary spectrum ($k<m/4$), a much stronger constraint.
\end{remark}

\subsection{Constant symbol matrix for circulant t(r)ienerments}
\label{ssec:symbol-matrix}

\begin{definition}[Circulant t(r)ienerment]
A t(r)ienerment $\tau$ on $\mathbb{Z}_n$ is \emph{circulant} if there exist disjoint
sets $S^+,S^-,S^0\subset\mathbb{Z}_n\setminus\{0\}$ with
$S^+\cup S^-\cup S^0=\mathbb{Z}_n\setminus\{0\}$ and $-S^+=S^-$ and $-S^0=S^0$,
such that:
\begin{itemize}
\item $i\to j$ iff $j-i\in S^+$,
\item $j\to i$ iff $j-i\in S^-$,
\item $i\leftrightarrow j$ iff $j-i\in S^0$ (tied pairs).
\end{itemize}
The group $\mathbb{Z}_n$ acts by rotation, giving the GLMY chain complex
$C_\bullet(D_\tau)$ a $\mathbb{Z}_n$-equivariant structure.
\end{definition}

\begin{theorem}[Constant symbol matrix for t(r)ienerments]\label{thm:symbol-const}
Let $\tau$ be a circulant t(r)ienerment on $\mathbb{Z}_n$.  Then the Tang--Yau
symbol matrix $M_m(t)$ of the associated digraph $D_\tau$ is constant
(independent of $t$), and hence all $n$ eigenspaces of the $\mathbb{Z}_n$-action
have identical dimensions $\dim\Omega_m^{(k)}$ for $k=0,\ldots,n-1$.

In particular, the Betti numbers satisfy:
\[
  \beta_p(D_\tau)=n\cdot\beta_p^{(0)}
\]
whenever all eigenspaces contribute equally to $\beta_p$
(verified for Paley $T_p$ at $p=7,11$; conjectured in general).
\end{theorem}

\begin{proof}
The proof is identical to that of THM-125 \cite{Cassidy2026} for circulant
tournaments.  The Tang--Yau symbol matrix is defined as
\[
  M_m[J,D](t)=\sum_{\substack{i:\text{face}_i(D)=J\\\text{face}_i \text{ non-degenerate}}}
              (-1)^i\,t^{\mathrm{offset}_i(D)}.
\]
The key lemma is: for any allowed path $D=(d_1,\ldots,d_m)\in A_m^{\tau}$,
the tail $(d_2,\ldots,d_m)$ is in $A_{m-1}^{\tau}$.

\smallskip
The allowed paths $A_m^{\tau}$ for the digraph $D_\tau$ consist of sequences
$(v_0,v_1,\ldots,v_m)$ of distinct vertices with each consecutive pair
$(v_i,v_{i+1})\in A(D_\tau)$.  For circulant $\tau$, $D_\tau$ is also circulant
(rotation by $\mathbb{Z}_n$ preserves arcs).  The argument of THM-125 applies
verbatim: the partial sums $\{s_j=\sum_{i=1}^j d_i\bmod n\}$ being all distinct
for $(d_1,\ldots,d_m)\in A_m^{\tau}$ implies the same for the tail.  This
argument uses only the injectivity of partial sums, which depends on the
distinctness of all vertices — a property of any path, regardless of whether
arcs are directed or tied (bidirectional arcs also define valid paths in $D_\tau$).

\smallskip
Therefore all non-trivial contributions to $M_m(t)$ come from inner face maps
(offset 0), making $M_m(t)$ constant.  The corollary about eigenspace dimensions
follows as in THM-125.
\end{proof}

\begin{corollary}[Betti speedup for circulant t(r)ienerments]\label{cor:speedup}
For any prime $n$ and circulant t(r)ienerment $\tau$ on $\mathbb{Z}_n$,
the Betti numbers of $D_\tau$ can be computed using a single eigenspace,
reducing computation by a factor of $n$.  Combined with sparse matrix techniques,
this gives the same asymptotic speedup as for circulant tournaments.
\end{corollary}

\subsection{The ternary metagraph as a group scheme quotient}
\label{ssec:scheme-quotient}

The tournament metagraph $G_n=Q_m^{(2)}/S_n$ is the quotient of the binary
$m$-cube by the symmetric group.  The t(r)ienerment metagraph
$G_n^{\mathrm{tri}}=Q_m^{(3)}/S_n$ is the quotient of the ternary $m$-cube.

\begin{proposition}[Structure of the ternary quotient map]
The quotient map $\pi:Q_m^{(3)}\to G_n^{\mathrm{tri}}$ is a $T(n)\cdot n!$-fold
covering in the generic case (when $|\mathrm{Aut}(\tau)|=1$).
The fiber size over $[\tau]$ is $n!/|\mathrm{Aut}(\tau)|$ (weighted by the
number of tilings, which equals $H_{\mathrm{dir}}(\tau)$ by
Theorem~\ref{thm:tiling-OS}).
\end{proposition}

The ternary cube $Q_m^{(3)}$ has richer structure than $Q_m^{(2)}$:
\begin{itemize}
\item It has a $\mathbb{Z}_3$-module structure (vs.\ $\mathbb{F}_2$-module for
  the binary cube), giving more symmetry.
\item The complement involution $t_j\mapsto 1-t_j$ of the binary cube extends to
  a $\mathbb{Z}_3$ rotation $t_j\mapsto t_j+1\pmod{3}$ that cyclically permutes
  states: $0\to 1\to 2\to 0$ (i.e., forward$\to$backward$\to$tie$\to$forward).
\item This $\mathbb{Z}_3$ rotation commutes with the $S_n$ action (since $S_n$
  permutes tile positions, not states), giving a $\mathbb{Z}_3$ action on
  $G_n^{\mathrm{tri}}$.
\end{itemize}

\begin{definition}[Ternary complement]
The \emph{ternary complement} of a t(r)ienerment $\tau$ is the t(r)ienerment
$\tau^{(3)}$ obtained by the cyclic state rotation: directed arcs become backward
arcs, backward arcs become ties, and ties become directed arcs.  Formally,
$t_j\mapsto t_j+1\pmod 3$ in the tiling model.
\end{definition}

\begin{proposition}[Ternary complement orbits]
Under repeated application of the ternary complement, every t(r)ienerment $\tau$
has a $\mathbb{Z}_3$-orbit of size 1 or 3.  Size-1 orbits are t(r)ienerments
$\tau\cong\tau^{(3)}$; these exist only when $n\equiv 0\pmod 3$.
\end{proposition}

\begin{proof}
Applying the ternary complement three times returns to the original state
($t_j+3\equiv t_j$), so every orbit has size dividing 3.  A size-1 orbit
requires $\tau^{(3)}\cong\tau$, meaning the state rotation is realized by
some vertex relabeling.  For this to be possible, the number of directed arcs
must equal the number of ties (since the rotation cyclically maps these sets).
Both count must be $\binom{n}{2}/3$, requiring $3\mid\binom{n}{2}=n(n-1)/2$,
i.e., $3\mid n(n-1)$, which holds for $n\equiv 0$ or $1\pmod 3$.
\end{proof}

\begin{remark}[Ternary merged metagraph]
The $\mathbb{Z}_3$-quotient $G_n^{\mathrm{tri}}/\mathbb{Z}_3$ is the natural
analog of the $\mathbb{Z}_2$-merged metagraph $G_n/\mathbb{Z}_2$ from
\cite{Cassidy2026}.  Its vertex count is
$V_{\mathrm{ter}}=(T(n)+2F_3(n))/3$
where $F_3(n)$ counts iso classes fixed by one application of the ternary
complement.  The merged metagraph has the spine-rib-sea decomposition
lifted to a $\mathbb{Z}_3$-equivariant setting.
\end{remark}

\subsection{Paley t(r)ienerments and quadratic residue structure}
\label{ssec:paley-trien}

\begin{definition}[Paley t(r)ienerment]\label{def:paley-trien}
For a prime $p\equiv 3\pmod 4$, the \emph{Paley t(r)ienerment} $\tilde{T}_p$
on $\mathbb{Z}_p$ is the circulant t(r)ienerment with:
\begin{itemize}
\item $i\to j$ (directed) iff $j-i$ is a quadratic residue mod $p$,
\item $i\leftrightarrow j$ (tied) iff $j-i$ is a quadratic non-residue mod $p$.
\end{itemize}
This is a circulant t(r)ienerment with $S^+=\mathrm{QR}$, $S^-=\emptyset$
(no backward directed arcs), $S^0=\mathrm{NQR}$ (non-residues are tied).
\end{definition}

\begin{remark}
The Paley t(r)ienerment $\tilde{T}_p$ has:
\begin{itemize}
\item $(p-1)/2$ directed arcs (the QR pairs), all in one direction,
\item $(p-1)/2$ tied pairs (the NQR pairs),
\item k = $(p-1)/2$ ties, so it lives in layer $G_p^{\mathrm{tri}}((p-1)/2)$.
\end{itemize}
The Paley tournament $T_p$ is NOT a sub-object of $\tilde{T}_p$: in $T_p$,
$i\to j$ for QR and $j\to i$ for NQR; in $\tilde{T}_p$, NQR pairs are tied.
\end{remark}

\begin{theorem}[Paley t(r)ienerment BIBD]\label{thm:paley-trien-bibd}
In the Paley t(r)ienerment $\tilde{T}_p$:
\begin{enumerate}
\item The directed 3-cycles of $D_{\tilde{T}_p}$ form a $2$-$(p,3,\tilde{\lambda}_3)$ BIBD.
\item The tie pairs $\{u,v\}$ form a $(p,(p-1)/2,(p-3)/4)$ partially balanced design.
\item $H(\tilde{T}_p)\ge H(T_p)$ (the t(r)ienerment has at least as many HPs as the
  Paley tournament).
\end{enumerate}
\end{theorem}

\begin{proof}[Proof of (1)]
In $D_{\tilde{T}_p}$, QR arcs are directed (one way) and NQR arcs are bidirectional.
A directed 3-cycle $(i,j,k)$ requires arcs $i\to j$, $j\to k$, $k\to i$ in $D_{\tilde{T}_p}$.
Each arc step can be: a QR step ($j-i\in\mathrm{QR}$, using directed arc) or an NQR step
($j-i\in\mathrm{NQR}$, using the $i\to j$ direction of the bidirectional arc).
The affine group $\mathrm{Aff}(\mathrm{QR})=\{x\mapsto ax+b:a\in\mathrm{QR},b\in\mathbb{Z}_p\}$
acts transitively on pairs, and the counting argument from THM-BIBD of \cite{Cassidy2026}
extends to $D_{\tilde{T}_p}$.  (The bidirectional NQR arcs contribute to directed cycles in
both orientations, increasing $\tilde{\lambda}_3>\lambda_3$ of the Paley tournament.)

(2): Tie pairs form a 2-design by the transitive group action on pairs.

(3): Every directed HP of $T_p$ remains a directed HP of $\tilde{T}_p$
(QR arcs are identical; NQR arcs become bidirectional, adding more traversal options).
\end{proof}

\begin{conjecture}[Paley t(r)ienerment maximizes $H$ in its layer]
Among all t(r)ienerments on $p$ vertices with $(p-1)/2$ ties (layer $k=(p-1)/2$),
the Paley t(r)ienerment $\tilde{T}_p$ maximizes $H$.
\end{conjecture}

This is the natural extension of the Paley maximization conjecture for tournaments.

\subsection{Summary of metagraph and eigenspace theorems}
\label{ssec:summary-eigenspace}

\begin{table}[ht]
\centering
\caption{Comparison of the binary (tournament) and ternary (t(r)ienerment) theories.}
\label{tab:comparison}
\begin{tabular}{p{5.2cm}p{5.2cm}}
\toprule
\textbf{Tournament theory (binary)} & \textbf{T(r)ienerment theory (ternary)}\\
\midrule
Tiling space $\{0,1\}^m=Q_m^{(2)}$ & Ternary tiling space $\{0,1,2\}^m=Q_m^{(3)}$\\
Orbit-stabilizer: $\#\text{tilings}\times|\mathrm{Aut}|=H$ & Ternary OS: $\#\text{tilings}\times|\mathrm{Aut}|=H_{\mathrm{dir}}$\\
Metagraph $G_n=Q_m^{(2)}/S_n$ & $G_n^{\mathrm{tri}}=Q_m^{(3)}/S_n$\\
$|V(G_n)|=\mathrm{A000568}(n)\sim 2^{\binom{n}{2}}/n!$ & $|V(G_n^{\mathrm{tri}})|=\mathrm{A007025}(n)\sim 3^{\binom{n}{2}}/n!$\\
Flip graph: binary Hamming $H(m,2)$ & Flip graph: ternary Hamming $H(m,3)$\\
Eigenvalues $m-2k$, multiplicity $\binom{m}{k}$ & Eigenvalues $2m-3k$, multiplicity $\binom{m}{k}2^k$\\
Spectral gap $2/m$; metagraph gap $\approx 2/n$ & Spectral gap $3/(2m)$; metagraph gap $\approx 3/(2n)$\\
Binary Krawtchouk $K_k^{(2)}$, $B_1\approx -H$ & Ternary Krawtchouk $K_k^{(3)}$, $B_1^{(3)}\approx -H$ (higher corr.)\\
Walsh degree $2\lfloor(n{-}1)/2\rfloor$ (THM-260) & Conjectured ternary degree $\lfloor(n{-}1)/2\rfloor$\\
$\mathbb{Z}_2$ complement, merged metagraph & $\mathbb{Z}_3$ ternary complement, $\mathbb{Z}_3$-merged metagraph\\
Circulant: constant symbol matrix (THM-125) & Circulant: same proof gives same result\\
$H(T)$ always odd (R\'edei) & $H(\tau)$ any positive integer\\
BIBD: 3-cycles form $2$-design in Paley $T_p$ & BIBD: extends to $D_{\tilde{T}_p}$ with more cycles\\
\bottomrule
\end{tabular}
\end{table}

\end{document}
